\documentclass[11pt]{article}
\usepackage[margin=0.9in]{geometry}
\usepackage{amsmath,amssymb,amsthm,booktabs}
\usepackage[T1]{fontenc}
\usepackage{lmodern,microtype}
\usepackage[colorlinks=true,linkcolor=blue,urlcolor=blue]{hyperref}
\newtheorem{theorem}{Theorem}[section]
\newtheorem{lemma}[theorem]{Lemma}
\newtheorem{proposition}[theorem]{Proposition}
\newtheorem{corollary}[theorem]{Corollary}
\theoremstyle{remark}
\newtheorem{remark}[theorem]{Remark}
\newcommand{\R}{\mathbb R}
\newcommand{\C}{\mathbb C}
\newcommand{\norm}[1]{\left\|#1\right\|}
\newcommand{\diag}{\operatorname{diag}}
\title{Toward a four-dimensional Yang--Mills construction\\
\large Notebook 07: locality and refinement limit of the flat Gaussian}
\author{Research working document}
\date{16 October 2026}
\begin{document}
\maketitle

\begin{abstract}
We derive uniform exponential spatial decay for the flat quadratic
effective action and its constrained Gaussian covariance. The action
is a gauge quotient of an averaged Feynman covariance. An exact
algebraic cancellation removes its apparent complex massless poles.
A discrete Parseval estimate controls the growing number of fine
Fourier aliases and supplies a complex momentum strip independent of
refinement. The fixed local gauge penalty and pivot insertion then
combine this locality with the spectral gap proved in Notebook 06.
We also prove convergence to an explicit continuum sinc-alias symbol
at rate $O(M^{-2})$ in an exponentially weighted kernel norm, uniformly
over compatible periodic volumes. The refinement and thermodynamic
limits commute for this free reference. These results supply the flat
baseline required by the earlier Gaussian cutoff theorems. They do not
control the actual complex-background perturbations or prove a nonlinear
continuum Yang--Mills limit.
\end{abstract}

\noindent\textbf{Status.}
Research proof checked by separate internal derivations and audits.
No external certification or novelty claim is made.
The flat variational model is defined explicitly
below. Its identification with the physical Gaussian retains the
smooth constrained-minimizer and coordinate requirements of Notebook 06;
Section~\ref{sec:source} records the precise deposited source match.

\section{The result and its scope}

Fix dimension $d$ and odd next-step blocking factor $L\ge3$.
Let $M=L^k$ be the refinement from physical spacing $\eta=M^{-1}$
to spacing one. More generally, any odd positive $M$ is allowed.
Use the straight-string average and physical volume norms of Notebook 06,
and define
\[
 q_M(B)=\inf_{Q_M^{\rm st}A=B}\norm{d_\eta A}_\eta^2.
\]
For finite periodic lattices this is a finite-dimensional quadratic
form. Its translation-invariant kernel and infinite-volume version
will follow from the Fourier calculation below.

Let $T_L$ be the averaged rooted-path map within fixed $L$-blocks,
$S_L=Q_L^\flat$ the next-step rooted constraint, and $K_L$ the
actual flat insertion of pivot bonds into the retained nonpivot
coordinates. The reduced precision is
\begin{equation}\label{eq:reduced}
 A_{M,0}=K_L^*[q_M+T_L^*T_L]K_L.
\end{equation}
All stars here are real Hilbert-space adjoints in the fixed volume
conventions. Their holomorphic continuations are bilinear transposes,
not variable-dependent complex conjugations.

\begin{theorem}\label{thm:flat-decay}
There are positive constants $\rho,C_q$, depending only on $d$, such that
the unit-lattice action kernel satisfies
\[
 \norm{q_M(x)}\le C_qe^{-\rho|x|_1}
\]
uniformly in $M$. There are positive $c,C$, depending only on fixed
$d,L$ and the fixed color norm, such that the retained-coordinate
covariance satisfies
\[
 \norm{A_{M,0}^{-1}(n,\alpha;m,\beta)}
 \le Ce^{-c|n-m|_1}.
\]
Here $n,m$ are $L$-block indices and $\alpha,\beta$ are retained
bond labels within a block. Finite periodic kernels are exact
periodizations, with uniform exponential bounds in periodic distance.
\end{theorem}

The theorem concerns the homogeneous full-block quadratic problem.
A boundary upper bound, such as that of Notebook 06, does not imply
that a boundary effective covariance is a restriction of this one.
An arbitrary retained-history or nonflat operator needs its own match.

\section{The exact alias formula and the gauge quotient}

Let $I_M=\{-(M-1)/2,\ldots,(M-1)/2\}$ and, for coarse momentum
$z$, write $p_r=z+2\pi r$, $r\in I_M^d$. Define the entire multipliers
\[
 a_M(p)=\prod_{\nu=1}^d\frac1M\sum_{s\in I_M}e^{ip_\nu s/M},
 \qquad
 \ell_{M,\mu}(p)=\frac1M\sum_{s=0}^{M-1}e^{ip_\mu s/M},
\]
\[
 \omega_\mu(p)=a_M(p)\ell_{M,\mu}(p).
\]
This is the actual one-form averaging multiplier. Its longitudinal
factor occurs in addition to the centered block factor in that
direction. A scalar block average alone would be the wrong operator.

Write
\[
 u_{r,\mu}^{\pm}(z)=M(e^{\pm ip_{r,\mu}/M}-1),\qquad
 P_r(z)=\sum_\mu u_{r,\mu}^+(z)u_{r,\mu}^-(z)
       =4M^2\sum_\mu\sin^2(p_{r,\mu}/(2M)),
\]
\[
 d_\mu^\pm(z)=e^{\pm iz_\mu}-1,\qquad
 b_\mu(z)=d_\mu^+(z)d_\mu^-(z)=4\sin^2(z_\mu/2).
\]
Minus means analytic substitution of the negative argument. It agrees
with conjugation on real momenta only. String telescoping gives
\begin{equation}\label{eq:gradient-symbol}
 \omega_\mu(p_r)u_{r,\mu}^+(z)=a_M(p_r)d_\mu^+(z),
\end{equation}
and the corresponding minus identity.

With physical-volume Parseval normalization,
$B_\mu(z)=\sum_r\omega_\mu(p_r)A_{r,\mu}$.
An embedded plane wave has the same fine-volume and coarse-volume
norm, so no extra $M^d$ factor enters the alias covariance.

For real $z\ne0$ in the central Brillouin zone, all $P_r>0$.
The fine Maxwell energy in one alias is
\[
 P_r|A_r|^2-| (u_r^+)^*A_r|^2
 =\inf_\lambda P_r|A_r+u_r^+\lambda|^2.
\]
By \eqref{eq:gradient-symbol}, fine gradients produce precisely coarse
gradients. They produce every scalar multiple of $d^+$ because the
central $a_M(z)$ is nonzero on this real zone. Thus Maxwell minimization
equals Feynman minimization with the coarse representative allowed to
vary by a gradient. The averaged Feynman covariance is diagonal:
\begin{equation}\label{eq:feynman}
 W(z)=\diag(c_\mu(z)),\qquad
 c_\mu(z)=\sum_{r\in I_M^d}
       \frac{\omega_\mu(p_r)\omega_\mu(-p_r)}{P_r(z)}.
\end{equation}
For real momenta its numerators are squared absolute values.
Minimizing the positive Feynman form at prescribed coarse datum $B'$
gives $B'^*W^{-1}B'$. A subsequent minimization over
$B'=B+d^+\lambda$ gives the full-space precision
\begin{equation}\label{eq:quotient}
 q_M=W^{-1}-W^{-1}d^+
       (d^-W^{-1}d^+)^{-1}d^-W^{-1}.
\end{equation}
This also follows by the ordinary scalar Schur complement.
Neither $W^{-1}$ nor the transverse compression of $W^{-1}$ alone
is generally this quotient.

At $z=0$, every nonzero alias has $a_M(2\pi r)=0$.
The average reads fine constants, which realize every constant
coarse field with zero Maxwell energy. Therefore
\begin{equation}\label{eq:zero-mode}
 q_M(0)=0.
\end{equation}
The harmonic sector is handled explicitly; no inverse of its zero
Feynman eigenvalue is assigned.

\section{Cancellation of the complex massless denominator}

Separate the central alias and define
\[
 A(z)=a_M(z)a_M(-z),\quad
 g_\mu(z)=\omega_\mu(z)\omega_\mu(-z),\quad
 R_\mu(z)=\sum_{r\ne0}
       \frac{\omega_\mu(p_r)\omega_\mu(-p_r)}{P_r(z)},
\]
\[
 D_\mu=g_\mu+P_0R_\mu,\qquad h_\mu=D_\mu^{-1}.
\]
The central gradient identity implies
\begin{equation}\label{eq:central-identity}
 \sum_\mu\frac{b_\mu}{g_\mu}=\frac{P_0}{A}.
\end{equation}
All removable one-variable factors are understood analytically.
Set
\begin{equation}\label{eq:E}
 E=A^{-1}-\sum_\mu\frac{b_\mu R_\mu}{g_\mu D_\mu}.
\end{equation}
Subtracting $D_\mu^{-1}-g_\mu^{-1}$ in
\eqref{eq:central-identity} proves exactly
\[
 \sum_\mu b_\mu h_\mu=P_0E.
\]
On the original real nonzero set, $W^{-1}=P_0\diag(h_\mu)$.
Substitution in \eqref{eq:quotient} and cancellation gives
\begin{equation}\label{eq:canceled}
 (q_M)_{\mu\nu}(z)
 =P_0(z)h_\mu(z)\delta_{\mu\nu}
  -\frac{h_\mu(z)d_\mu^+(z)d_\nu^-(z)h_\nu(z)}{E(z)}.
\end{equation}
There is no inverse power of $P_0$ left. This removes its entire
complex null set, rather than just the real zero at the origin.
We will prove that the remaining denominators are uniformly nonzero
on a complex strip. The resulting formula is zero at $z=0$, agreeing
with \eqref{eq:zero-mode}. For $M=1$, $A=g_\mu=D_\mu=E=1$
and $R_\mu=0$, giving the usual finite-range Maxwell precision.

\section{A complex momentum strip independent of refinement}

\subsection{Alias Parseval prevents a factor $M^d$}

Discrete orthogonality gives, for real $x,y$ in one coordinate,
\[
 \sum_{r\in I_M}\left|
       \frac1M\sum_{s\in I_M}e^{i(x+2\pi r+iy)s/M}
                         \right|^2
 =\frac1M\sum_{s\in I_M}e^{-2ys/M}\le e^{|y|}.
\]
Taking products and using $|\ell_{M,\mu}(x+2\pi r+iy)|\le
e^{|y_\mu|}$ proves
\begin{equation}\label{eq:alias-parseval}
 \sum_r|\omega_\mu(x+2\pi r+iy)|^2
 \le e^{\sum_\nu|y_\nu|+2|y_\mu|}.
\end{equation}
The same bound holds for the negative arguments. Cauchy--Schwarz
therefore bounds the sum of absolute products in \eqref{eq:feynman}
by the same expression. The number of aliases does not appear.

\subsection{High aliases and real denominator gaps}

Choose $K=\pi+\varepsilon$ with $0<\varepsilon\le1/4$.
For real $|x_\mu|\le K$, every nonzero alias has a coordinate whose
distance from $2\pi M\mathbb Z$ is at least $2\pi-K$. Hence
\[
 P_r(x)\ge H_K:=4(2\pi-K)^2/\pi^2>0\quad(r\ne0).
\]
For complex perturbations,
\begin{align*}
 \Re P_r(x+iy)
 &=P_r(x)-2M^2\sum_\mu\cos((x_\mu+2\pi r_\mu)/M)
                                  [\cosh(y_\mu/M)-1]\\
 &\ge P_r(x)-e^{\delta}\sum_\mu y_\mu^2
 \qquad(|y_\mu|\le\delta\le1).
\end{align*}
Choose a positive $\delta$ so $de^\delta\delta^2<H_K/2$.
All high-alias denominators then have real part at least $H_K/2$.
Equation \eqref{eq:alias-parseval} proves that each $R_\mu$ is
holomorphic and uniformly bounded on this open tube. The central
finite averages and $P_0$ are uniformly bounded there too. For example
$|M(e^{iz/M}-1)|\le |z|e^{|\Im z|/M}$.

On the enlarged real cube, set
\[
 c_K=\frac{\sin(K/2)}{K/2}>0,\qquad
 g_*=c_K^{2d+2},\qquad D_*=1+dK^2/H_K.
\]
The centered factor
$r_M(t)=\sin(t/2)/[M\sin(t/(2M))]$, with $r_M(0)=1$,
satisfies $c_K\le r_M(t)\le1$ for $|t|\le K$.
Using real alias Parseval gives
\begin{equation}\label{eq:real-gaps}
 A\ge c_K^{2d},\qquad
 g_\mu,D_\mu\ge g_*,\qquad
 0\le R_\mu\le H_K^{-1},\qquad D_\mu\le D_*.
\end{equation}
For real nonzero $x$,
\[
 E(x)=\frac{\sum_\mu b_\mu(x)/D_\mu(x)}{P_0(x)}
      \ge\frac{c_K^2}{D_*},
\]
since $|d_\mu^+(x)|\ge c_K|x_\mu|$ and $P_0(x)\le|x|^2$.
At zero, \eqref{eq:E} gives $E(0)=1$.

\subsection{Two successive shrinkings}

Take nested real cubes strictly between $[-\pi,\pi]^d$ and
$(-K,K)^d$, with nested positive imaginary widths. Cauchy's estimates
on the initial tube give uniform derivatives for $A,g_\mu,D_\mu$
on the first smaller tube. Integrating along $x+ity$ and using
\eqref{eq:real-gaps} keeps these functions uniformly away from zero
on a still smaller tube, independently of $M$.

Formula \eqref{eq:E} is now holomorphic and uniformly bounded.
Apply Cauchy's estimate once more, using its real lower bound, to keep
$E$ uniformly away from zero on a final tube. Its real cube still
strictly contains $[-\pi,\pi]^d$. Thus some fixed $\rho_0>0$ gives
a uniformly bounded holomorphic formula \eqref{eq:canceled} for
$|\Im z_\mu|<\rho_0$, with a positive real margin around the central
Brillouin cube. The order of these steps avoids assuming the desired
nonvanishing of $E$ while estimating it. No Cauchy disk leaves an
already controlled open domain.

\section{Periodic gluing, kernels, and the reduced inverse}

The central-alias split is not individually $2\pi$ periodic.
The full quotient is periodic on real momenta: shifting by $2\pi$
permutes alias labels modulo $M$, the fine symbols have period
$2\pi M$, and the coarse gradient is unchanged. The continuous
zero-mode value agrees as well. The analytic extensions on adjacent
enlarged Brillouin charts consequently agree on their open real
overlaps. Applying the identity theorem successively in the complex
coordinates makes them agree holomorphically. The positive real margin
is what supplies genuine overlaps. They therefore define one bounded
holomorphic periodic symbol on a strip of fixed width.

For its infinite-lattice coefficient
\[
 q_M(x)=(2\pi)^{-d}\int_{[-\pi,\pi]^d}e^{iz\cdot x}q_M(z)\,dz,
\]
shift each contour in the sign of $x_\mu$. For every fixed
$0<\rho<\rho_0$ this gives
\[
 \norm{q_M(x)}\le C_qe^{-\rho|x|_1}.
\]
Reducing the exponent gives an absolutely summable weighted kernel.
Finite tori evaluate the same variational symbol on their momentum
grids, so their kernels are exactly its periodizations. Absolute
convergence justifies every image sum. The weighted row sum on a
torus is bounded by the infinite weighted row sum because periodic
distance is at most the distance of each representative. These
bounds include short periods with the correct oriented-cell
multiplicities.

\subsection{Direct infinite-volume variational identification}

The symbol also gives the actual infinite-volume infimum, without
interchanging an infimum and a volume limit. For real nonzero momentum
let $h=\diag(h_\mu)$ and
\[
 \theta=\frac{d^-hB}{d^-hd^+},\qquad B'=B-d^+\theta.
\]
The real upper and lower bounds on $D_\mu$ imply
$|d^+\theta|\le C|B|$ and $|B'|\le C|B|$, uniformly in $M$.
The Feynman minimizing fiber for $B'$ is
\[
 A_r^F=\diag(\omega_\mu(-p_r))\frac{P_0}{P_r}hB'.
\]
For the central alias the ratio is one; for the others it is bounded
by $d\pi^2/4$ on the central real cube. Alias Parseval bounds
$\sum_r|A_r^F|^2\le C|B|^2$.
Moreover $d^-hB'=0$ and \eqref{eq:gradient-symbol} imply that every
$A_r^F$ is fine transverse. Its Maxwell and Feynman energies agree.

Restore the exact coarse datum by adding only the central fine gradient:
\[
 A_0=A_0^F+u_0^+\theta/a_M(z),\qquad A_r=A_r^F\ (r\ne0).
\]
Its average is $B$ and its Maxwell energy is unchanged. The added
bond field is uniformly bounded because
$|u_0^+|/|d^+|\le\pi/2$ and $|a_M(z)|\ge(2/\pi)^d$.
The scalar $\theta$ can be singular at zero, but this actual bond
field is not. These measurable minimizing fibers define a fine
$L^2$ field for every coarse $L^2$ field. The zero point has measure
zero in infinite volume; on a finite torus use its constant lift.
This proves the required variational identification directly.

\subsection{The fixed local maps introduce no refinement dependence}

$T_L$ uses paths inside one $L$-block. Each $S_L$ row uses only
the two endpoint blocks of its coarse bond. Its center-line interface
pivot has discrete coefficient $L^{-d}$, and no other row uses that
pivot. Thus $K_L$ is the identity on retained coordinates and has
pivot rows $-L^d(S_L)_{\rm free}$. It is finite range, with coefficients
depending only on fixed $d,L$.

Group retained coordinates by the $L$-block containing their bond
basepoint. There are $d(L^d-1)$ scalar retained coordinates per block
in this convention, with a fixed extra color multiplicity. Since
unit-site distance between blocks $n,m$ is at least
$L|n-m|_1-d(L-1)$, the action kernel and these fixed finite-range maps
give constants $\eta_0,C_A>0$, independent of $M$, with
\begin{equation}\label{eq:precision-schur}
 \sum_n e^{\eta_0|n|_1}\norm{A_{M,0}(n)}\le C_A.
\end{equation}
This is the locality input that the spectral gap alone did not supply.

Notebook 06 proves $A_{M,0}\ge\gamma_*I$ uniformly in refinement
and volume. Hence the finite-dimensional block Bloch symbol obeys
$A_M(\kappa)\ge\gamma_*I$ for every real $\kappa$.
One can infer this either from the infinite-volume form inequality
or from its values on the dense union of finite-torus momentum grids
and continuity of the symbol.

By \eqref{eq:precision-schur}, $A_M(\kappa)$ is holomorphic for
$|\Im\kappa|_\infty<\eta_0$. For $|y|_\infty\le\eta_0/2$,
\[
 \norm{A_M(x+iy)-A_M(x)}
 \le |y|_\infty\sum_n|n|_1 e^{(\eta_0/2)|n|_1}
                                      \norm{A_{M,0}(n)}
 \le C_1|y|_\infty,
\]
where $C_1$ is uniform. Choose a positive
$\eta_C<\min(\eta_0/2,\gamma_*/(2C_1))$, omitting the second
restriction if $C_1=0$. A Neumann inverse about the real symbol gives
\[
 A_M(\kappa)^{-1}\text{ holomorphic},\qquad
 \norm{A_M(\kappa)^{-1}}\le2/\gamma_*
 \quad(|\Im\kappa|_\infty<\eta_C).
\]
Contour shifting proves exponential decay of its block kernel and
completes Theorem~\ref{thm:flat-decay}.

Finite tori sample this same inverse symbol. Their covariance is the
periodization of the infinite inverse kernel, rather than the inverse
of a truncated infinite covariance. On each fixed region the image
error tends to zero exponentially as every period tends to infinity,
uniformly in $M$ when comparing the same refinement. No large-period
assumption was needed for the uniform bound itself.

\section{The free low-momentum normalization}

Let $U_\mu(z)=e^{iz_\mu/2}$ and pass to the bond-midpoint convention
$q_M^{\rm mid}(z)=U(-z)q_M(z)U(z)$. The canceled formula becomes
\[
 (q_M^{\rm mid})_{\mu\nu}(z)
 =P_0h_\mu\delta_{\mu\nu}
  -\frac{4h_\mu\sin(z_\mu/2)\sin(z_\nu/2)h_\nu}{E}.
\]
The functions $g_\mu,R_\mu$ are even in every momentum coordinate;
for $R_\mu$, reflect the corresponding alias index in its symmetric
set. The phases cancel in each product of positive and negative
arguments. Thus $h_\mu,E$ are even too, and
$h_\mu(0)=E(0)=1$, $R_\mu(0)=0$.
Uniform Cauchy bounds in the established neighborhood give
\[
 h_\mu(z)=1+O(|z|^2),\qquad E(z)=1+O(|z|^2),\qquad
 P_0(z)=\sum_\mu z_\mu^2+O(|z|^4),
\]
with constants independent of $M$. Hence, for small real momentum,
\[
 q_M^{\rm mid}(p)=|p|^2I-pp^T+O(|p|^4).
\]
This checks the Maxwell normalization of the flat quadratic model.
It is not a statement about an interacting renormalized coupling
or the nonlinear continuum limit.

\section{The actual flat source covariance}\label{sec:source}

The maps must remain distinct. $Q_M^{\rm st}$ goes from original
spacing $M^{-1}$ to unit spacing; the rooted composition has the
same energy infimum through the explicit gradient adjustment of
Notebook 06. $S_L,T_L,K_L$ concern the next unit-to-$L$ step.
The averaged Feynman covariance $W$ is an intermediate calculation,
not the actual reduced fluctuation covariance.

Raw II, the small-field step around lines 73860--74020, records the
quadratic expansion, exponential gauge fixing, straightening of the
nonlinear constraint, coupling rescaling, and pivot elimination of
CMP109 (2.11)--(2.13). At $(U,J)=(1,0)$, the stationary action
Hessian is $q_M$, the gauge Hessian is $T_L^*T_L$, and the constraint
straightening starts at quadratic order. Its second derivative
therefore does not change the stationary Hessian. Haar and Jacobian
terms remain in the fluctuation interaction. The retained Gaussian
coordinates are exactly the nonpivot variables inserted by $K_L$.
Under the smooth constrained-minimizer identification of Notebook 06,
the source covariance is consequently
\[
 C^{(k)}(1)=\bigl(K_L^*[q_M+T_L^*T_L]K_L\bigr)^{-1}.
\]

The kernel convention matters: Raw I, lines 46960--47000, writes
the discrete next-step matrix coefficient as $L^{-d}Q^\flat(c,b)$.
At flatness its pivot kernel is the identity, so the actual insertion
coefficient is $-Q^\flat(c,b)=-L^dS_L(c,b)$ on nonpivot columns,
exactly as used here. No nonlocal inverse is hidden in this pivot map.

The source's differently gauge-fixed auxiliary $G_1$ is not the
Feynman inverse used in \eqref{eq:feynman}. Its transcribed identities
$G_1^{-1}=\Delta_1+DRD^*+Q^*aQ$,
$H_1=G_1Q^*(QG_1Q^*)^{-1}$, $QH_1=I$, and $RD^*H_1=0$
instead give $H_1^*\Delta_1H_1=(QG_1Q^*)^{-1}-a$.
The variational source match identifies that physical action Hessian
with $q_M$; it does not identify the two auxiliary propagators.

The result supplies a uniformly positive, bounded, exponentially
local flat covariance on the exact retained coordinates of this
full-block step. What remains for the Gaussian cutoff application
is weighted control of the actual complex independent-current
precision perturbation or its source responses. The flat proof
does not provide those derivatives, arbitrary-reference background
locality, mixed-history domain bounds, or permission to change a
sharp cutoff's integration coordinates. Overlapping interactions
and the continuum Yang--Mills obligations also remain separate.

\begin{thebibliography}{9}
\bibitem{repo} \href{https://github.com/ethanwsussman/Yang--Mills-LLM-experiment/tree/e392ef45013a581664bab76b90867170b30fe10e}
{Pinned Yang--Mills-LLM-experiment account}, commit
\texttt{e392ef45013a581664bab76b90867170b30fe10e}.
\bibitem{balaban109} T. Ba\l aban,
\emph{Renormalization group approach to lattice gauge field theories I},
Commun. Math. Phys. 109 (1987), 249--301.
\href{https://scispace.com/pdf/renormalization-group-approach-to-lattice-gauge-field-3ygai0vxt1.pdf}
{Copy of the original article}.
\end{thebibliography}
\end{document}
